\documentclass[italian,12pt]{article}

\def\Title{Verbale interno}
\def\Subtitle{Riunione interna settimanale}
\def\Author{7Last}
\def\Date{2024-05-22}
\def\Version{v1.0}

%pacchetti extra da scaricare dblfloatfix, fancyhdr
\usepackage[left=2cm, right=2cm, bottom=3cm, top=3cm]{geometry}
\usepackage{fancyhdr}%creazione header-footer
\usepackage{graphicx} %serve per inserire immagini
\graphicspath{ {../../../logo/} }
%\usepackage{dblfloatfix} %serve per posizionare gli elementi dove si vuole
\usepackage[hidelinks]{hyperref} %serve per i link
\usepackage{tikz}
\usepackage{tgadventor} % font
\usepackage[useregional=numeric,showseconds=true,showzone=false]{datetime2}
\usepackage{caption}

\usepackage{hyperref}
\usepackage{tocloft}
\usepackage{titlesec}
\usepackage{color}
\usepackage{ulem}
\usepackage{pgfplots}
\usepackage{pgf-pie}
\usepackage[italian]{babel}
\usepackage{comment}
\usepackage{tabularx}
\usepackage{longtable}
\usepackage{float}
\usepackage{amsmath}

\linespread{1.2}
\captionsetup[table]{labelformat=empty}
\geometry{headsep=1.5cm}

\renewcommand{\contentsname}{Indice}
\renewcommand\familydefault{\sfdefault}


\glstoctrue

\glsaddkey
{feminine}
{}
{\glsentryfem}
{\Glsentryfem}
{\glsf}
{\Glsf}
{\GLSf}


\glsaddkey
{feminine_plural}
{}
{\glsentryfemplur}
{\Glsentryfemplur}
{\glsfp}
{\Glsfp}
{\GLSfp}



\setglossarysection{subsubsection}
\renewcommand{\glossarysection}[2][]{}

\makeglossaries


\newglossaryentry{ClickUp}{
    name=ClickUp,
    description={Software di gestione dei progetti che offre varie funzionalità, tra cui la gestione di task, una lavagna virtuale, fogli di calcolo e strumenti collaborativi per la creazione e modifica di documenti, il tutto accessibile da una piattaforma unificata.},
}

\newglossaryentry{Glossario}{
    name=Glossario,
    description={Elenco strutturato di termini tecnici o specializzati, ognuno corredato dalla propria definizione o spiegazione. Questo strumento aiuta a migliorare la comunicazione tra le varie parti coinvolte in un progetto, riducendo le ambiguità e garantendo una comprensione condivisa dei termini utilizzati in un determinato contesto.}
}

\newglossaryentry{Norme di Progetto}{
    name={Norme di Progetto},
    description={Regole e linee guida stabilite all'interno di un progetto per garantire coerenza e qualità nelle attività svolte. Definiscono standard e procedure, come documentazione, gestione delle versioni e criteri di codifica, per assicurare uniformità nell'approccio e nel risultato finale.}
}


\newglossaryentry{Piano di Progetto}{
    name={Piano di Progetto},
    description={Documento formale che delinea in dettaglio la pianificazione, la esecuzione, il monitoraggio e il controllo di tutte le attività coinvolte nella realizzazione di un progetto. Questo documento fornisca una roadmap chiara e organizzata, comprensiva di obiettivi, risorse, scadenze e strategie di gestione dei rischi. Essenziale per la gestione efficace di un progetto, il piano di progetto serve come guida per il team di lavoro e gli stakeholder, fornendo una struttura che facilita il coordinamento delle attivitå e l'assegnazione delle risorse.}
}

\newglossaryentry{Piano di Qualifica}{
    name={Piano di Qualifica},
    description={Documento che stabilisce gli standard di qualità, i processi e le attività di testing che saranno implementati durante lo sviluppo di un progetto. Contiene una descrizione dettagliata delle strategie di testing, delle metriche di valutazione e dei criteri di accettazione del prodotto finale. L'obiettivo principale del Piano di Qualifica è garantire che il prodotto soddisfi gli standard di qualità prefissati e che il processo di sviluppo segua procedure coerenti e efficaci.}
}


\newglossarystyle{myaltlistgroup}{
    \setglossarystyle{altlistgroup}
    \renewcommand*{\glsgroupheading}[1]{
        \pagebreak
        \item\makebox[\linewidth]{\Large\textbf{\glsgetgrouptitle{##1}}}
        \vspace*{-\baselineskip}
        \item\makebox[\linewidth]{\hspace*{3cm}\hrulefill\hspace*{3cm}}
        \addcontentsline{toc}{section}{\glsgetgrouptitle{##1}}
    }
}

\glsaddallunused


\begin{document}

    \newgeometry{left=2cm,right=2cm,bottom=2.1cm,top=2.1cm}
\begin{titlepage}
    \vspace*{.5cm}

    \vspace{2cm}
    {
        \centering
        {\bfseries\huge \Title\par}
        \bigbreak
        {\bfseries\large \Author\par}
        \bigbreak
        {\Version\par}
        \vfill

        \begin{tikzpicture}[remember picture,overlay]

            \fill[blue!20!black] (current page.south west) -- ++(10cm,0) -- ++(-10cm,15cm) -- cycle;
            \fill[orange] (current page.south east) -- ++(-18cm,0) -- ++(21.6cm,22cm) -- cycle;

            \clip (0,-3cm) circle (2.5cm) node (current page.center) {\includegraphics[width=5cm]{logo.jpg}};
        \end{tikzpicture}

    }

    \vfill

\end{titlepage}

\restoregeometry






















    \newpage

    %---------------header------------------%
\pagestyle{fancy}
\fancyhead{} % pulizia degli header
\lhead{%
\begin{tikzpicture}
    \clip (0,0) circle (0.5cm);
    \node at (0,0) {\includegraphics[width=1cm]{logo}};
\end{tikzpicture}%
}
\chead{\vspace{\fill}\Title\vspace{\fill}}
\rhead{\vspace{\fill}\Version\vspace{\fill}}
%quad è una spaziatura
%---------------------------------------%

    \begin{table}[!h]
        \caption{Versioni}
        \begin{center}
            \begin{tabular}{ c c c p{9cm} }
                \hline                                                                                               \\[-2ex]
                Ver. & Data       & Autore         & Descrizione                                                     \\
                \\[-2ex] \hline \\[-1.5ex]
                0.2  & 2024-04-02 & Elena Ferro    & Modifica stile documento, automatizzazione creazione glossario. \\
                0.1  & 2024-03-28 & Leonardo Baldo & Creazione documento                                             \\
                \\[-1.5ex] \hline
            \end{tabular}
        \end{center}
    \end{table}

    \newpage

    \tableofcontents

    \printglossary[style=myaltlistgroup,title=]
\end{document}






